% !TeX program = pdflatex
% Umowa akcjonariuszy Basilisk Systems P.S.A. Odesłania do umowy spółki Basilisk Systems (psa.tex, makro \US).
\documentclass[12pt,a4paper]{report}
\usepackage[T1]{fontenc}
\usepackage[utf8]{inputenc}
\usepackage[english,polish]{babel}
\usepackage{basilisk_i18n}
\usepackage{lmodern}
\usepackage{microtype}
\usepackage{setspace}
\usepackage{parskip}
\usepackage{needspace}
\usepackage{tabularx}
\usepackage[normalem]{ulem}

\BusinessPlanCompany{Basilisk Systems}
\BusinessPlanWatermarkText{Basilisk Systems}
\BusinessPlanWatermarkColor{black!6}
\BusinessPlanWatermarkFontSize{38}{46}
\BusinessPlanWatermarkAngle{38}
\BusinessPlanAuthorsText{Basilisk Systems --- projekt umowy akcjonariuszy}
\BusinessPlanConfidentialText{Projekt umowy akcjonariuszy --- poufne}
\BusinessPlanFooterText{Basilisk Systems --- umowa akcjonariuszy (projekt)}
\BusinessPlanVersionText{Wersja 1.7 --- projekt}

\setstretch{1.18}
\setlist[itemize]{topsep=0.2em,itemsep=0.15em,parsep=0.1em,leftmargin=2.0em}
\setlist[enumerate]{topsep=0.25em,itemsep=0.18em,parsep=0.1em,leftmargin=2.2em}
\setlength{\parindent}{0pt}
\setlength{\parskip}{0.55em}
\setlength{\emergencystretch}{4em}
\renewcommand{\arraystretch}{1.18}

\newcounter{paragraf}
\newcommand{\paragraf}[1]{%
  \refstepcounter{paragraf}%
  \phantomsection%
  \addcontentsline{toc}{section}{\S~\theparagraf\ #1}%
  \Needspace{6\baselineskip}%
  \subsection*{\S~\theparagraf\ #1}%
}
\newenvironment{ContractClause}[1]{%
  \paragraf{#1}%
}{%
  \par
}

\newcommand{\field}[1]{\textbf{[#1]}}
\newcommand{\CompanyFull}{Basilisk Systems prosta spółka akcyjna}
\newcommand{\CompanyShort}{Basilisk Systems P.S.A.}
% Odesłanie do paragrafu umowy spółki (numeracja psa.tex)
\newcommand{\US}[1]{§~#1 Umowy Spółki}
\newcommand{\USen}[1]{\S~#1 of the Articles}

\hypersetup{
  pdftitle={Basilisk Systems P.S.A. --- projekt umowy akcjonariuszy},
  pdfauthor={Basilisk Systems --- projekt umowy akcjonariuszy}
}

\begin{document}

\BusinessPlanTitlePage
  {Basilisk Systems}
  {Wersja 1.7 --- projekt}
  {Umowa akcjonariuszy}
  {Projekt umowy akcjonariuszy Basilisk Systems P.S.A.}
  {Zaangażowanie, vesting, własność intelektualna, poufność}
  {1.7}

\tableofcontents
\newpage

\chapter*{Projekt umowy akcjonariuszy}
\addcontentsline{toc}{chapter}{Projekt umowy akcjonariuszy}

\begin{center}
{\Large\bfseries UMOWA AKCJONARIUSZY}\par
\vspace{0.25cm}
{\LARGE\bfseries \CompanyFull}\par
\vspace{0.6cm}
\end{center}

zawarta w dniu \field{data} w \field{miejscowość} pomiędzy:

\begin{center}
\begin{tabularx}{\textwidth}{>{\bfseries}p{1.6cm} >{\raggedright\arraybackslash}X >{\raggedright\arraybackslash}p{3.2cm}}
\toprule
Strona & Dane identyfikacyjne & Akcje serii A \\
\midrule
A & \field{imię, nazwisko, PESEL, adres, adres e-mail} & \field{liczba, numery} \\
B & \field{imię, nazwisko, PESEL, adres, adres e-mail} & \field{liczba, numery} \\
C & \field{imię, nazwisko, PESEL, adres, adres e-mail} & \field{liczba, numery} \\
D & \field{imię, nazwisko, PESEL, adres, adres e-mail} & \field{liczba, numery} \\
E & \field{imię, nazwisko, PESEL, adres, adres e-mail} & \field{liczba, numery} \\
\bottomrule
\end{tabularx}
\end{center}

zwanymi dalej łącznie „Założycielami”, a każdy z osobna „Założycielem”, oraz

\textbf{\CompanyFull} z siedzibą w Pruszkowie, wpisaną do rejestru przedsiębiorców Krajowego Rejestru Sądowego pod numerem KRS \field{numer} (do czasu wpisu --- jako spółka w organizacji), reprezentowaną przez \field{imię, nazwisko} --- pełnomocnika powołanego uchwałą akcjonariuszy z dnia \field{data}, objętą aktem zawiązania Spółki (art.~300$^{79}$ §~1 KSH, \US{11} ust.~3 i \US{17} ust.~6), niebędącego Stroną niniejszej umowy i umocowanego do jej zawarcia wraz z Załącznikami A--D, zwaną dalej „Spółką”.

Strony zawierają umowę, której celem jest uporządkowanie współpracy Założycieli w sprawach, które nie należą do umowy spółki: zakresu zaangażowania, wykonania obowiązku zbycia akcji, własności intelektualnej, poufności i zakazu konkurencji. Umowa nie wprowadza praw uprzywilejowanych, prawa przyłączenia, przymusowego współzbycia ani innych mechanizmów inwestorskich.

\begin{ContractClause}{Definicje i stosunek do Umowy Spółki}
\label{sha:definicje}
\begin{enumerate}
  \item „Umowa Spółki” oznacza umowę \CompanyFull{} w brzmieniu obowiązującym w danej chwili. Terminy pisane wielką literą, niezdefiniowane w niniejszej umowie, mają znaczenie nadane im w Umowie Spółki, w szczególności: Założyciel Pierwotny, Założyciel Funkcjonalny Serii~F, Akcjonariusz Funkcjonalny, Dzień Przyznania Serii~F, Rada Dyrektorów i Walne Zgromadzenie.
  \item „Akcje Objęte” oznaczają akcje Akcjonariusza Funkcjonalnego objęte obowiązkiem zbycia na podstawie \US{9}, a „Akcje Niezwolnione” --- Akcje Objęte, które w danej chwili nie zostały zwolnione z tego obowiązku zgodnie z \US{9} ust.~2 albo~5.
  \item „Dzień Ustania Zaangażowania” oznacza dzień, w którym ustało stałe, osobiste zaangażowanie operacyjne Akcjonariusza Funkcjonalnego w rozwój Spółki, stwierdzony zgodnie z §~\ref{sha:ustanie}.
  \item Niniejsza umowa wiąże Strony jako zobowiązanie wzajemne. Nie zmienia Umowy Spółki i nie wiąże osób, które do niej nie przystąpiły. W razie sprzeczności między niniejszą umową a Umową Spółki w stosunkach ze Spółką rozstrzyga Umowa Spółki, a Strony zobowiązują się głosować za zmianą Umowy Spółki usuwającą sprzeczność w sposób najbliższy celowi niniejszej umowy.
  \item Niniejsza umowa stanowi umowę wykonawczą, o której mowa w \US{9} ust.~6. Nie jest umową przeniesienia praw własności intelektualnej; prawa te przechodzą na Spółkę wyłącznie na podstawie odrębnych umów, o których mowa w §~\ref{sha:ip}.
\end{enumerate}
\end{ContractClause}

\begin{ContractClause}{Przystąpienie nowych akcjonariuszy}
\label{sha:przystapienie}
\begin{enumerate}
  \item Założyciel zbywa akcje wyłącznie osobie, która przed zbyciem przystąpi do niniejszej umowy przez złożenie oświadczenia według wzoru z Załącznika C. Obowiązek ten nie dotyczy zbycia w postępowaniu egzekucyjnym ani nabywcy, wobec którego Walne Zgromadzenie, większością 75\% wszystkich głosów, zwolni z niego zbywcę.
  \item Założyciel Funkcjonalny Serii~F oraz osoba obejmująca akcje serii~F przystępuje do niniejszej umowy najpóźniej w dniu zawarcia umowy objęcia akcji. Do uczestników programu motywacyjnego (seria~P) niniejszej umowy nie stosuje się; ich obowiązki określa regulamin programu i umowa uczestnictwa. Uczestnik programu, który nabywa Akcje Objęte, przystępuje do niniejszej umowy zgodnie z ust.~1 i~3.
  \item Przystąpienie następuje w całości, przez oświadczenie według wzoru z Załącznika C przyjęte w imieniu wszystkich Stron przez Spółkę, którą Strony niniejszym do tego upoważniają; przystępujący nabywa prawa i obowiązki Strony, a zbywca nie zostaje przez to zwolniony z obowiązków powstałych przed zbyciem. Nabywca Akcji Objętych przejmuje obowiązki związane z tymi akcjami w zakresie, w jakim obciążały zbywcę, w tym składa ofertę i udziela pełnomocnictwa według wzorów z Załączników A i~B.
\end{enumerate}
\end{ContractClause}

\begin{ContractClause}{Zaangażowanie Założycieli}
\label{sha:zaangazowanie}
\begin{enumerate}
  \item Założyciele wskazani w Załączniku nr~1 do Umowy Spółki jako zobowiązani do stałego, osobistego zaangażowania operacyjnego obejmują następujące role i poświęcają Spółce co najmniej następujący czas:
\end{enumerate}

\begin{center}
\begin{tabularx}{\textwidth}{>{\bfseries}p{2.1cm} >{\raggedright\arraybackslash}X >{\raggedright\arraybackslash}p{3.2cm} >{\raggedright\arraybackslash}p{2.9cm}}
\toprule
Założyciel & Rola i obszar odpowiedzialności & Minimalny zakres & Podstawa prawna \\
\midrule
A & \field{rola} & \field{godz./tydz. albo etat} & \field{umowa o pracę / B2B / powołanie} \\
B & \field{rola} & \field{godz./tydz. albo etat} & \field{umowa o pracę / B2B / powołanie} \\
C & \field{rola} & \field{godz./tydz. albo etat} & \field{umowa o pracę / B2B / powołanie} \\
D & \field{rola} & \field{godz./tydz. albo etat} & \field{umowa o pracę / B2B / powołanie} \\
E & \field{rola} & \field{godz./tydz. albo etat} & \field{umowa o pracę / B2B / powołanie} \\
\bottomrule
\end{tabularx}
\end{center}

\begin{enumerate}[resume]
  \item Zaangażowanie jest wykonywane na podstawie umowy o pracę, umowy o świadczenie usług albo pełnienia funkcji dyrektora. Wynagrodzenie Założycieli ustala Rada Dyrektorów w budżecie, a wynagrodzenie dyrektorów --- uchwała akcjonariuszy. Do czasu pozyskania finansowania Założyciel może pełnić funkcję dyrektora bez wynagrodzenia. Niniejsza umowa nie wyłącza ani nie ogranicza wynagrodzenia należnego ze stosunku pracy ani wynikającego z bezwzględnie obowiązujących przepisów o świadczeniu usług; obniżenie wynagrodzenia wymaga zmiany właściwej umowy i nie działa wstecz.
  \item Zmiana roli albo zakresu zaangażowania wymaga zgody Rady Dyrektorów i zainteresowanego Założyciela. Obniżenie zakresu poniżej połowy minimum z ust.~1 na okres dłuższy niż 3 miesiące, bez zgody Rady Dyrektorów, stanowi ustanie zaangażowania.
  \item Inna działalność zawodowa Założyciela jest dopuszczalna, jeżeli nie narusza minimalnego zakresu zaangażowania, zakazu konkurencji (\US{14} ust.~5) ani obowiązku poufności.
\end{enumerate}
\end{ContractClause}

\begin{ContractClause}{Ustanie zaangażowania i rodzaj odejścia}
\label{sha:ustanie}
\begin{enumerate}
  \item Zaangażowanie Akcjonariusza Funkcjonalnego ustaje w szczególności z dniem:
  \begin{enumerate}[label=\alph*)]
    \item rozwiązania albo wygaśnięcia stosunku prawnego, na podstawie którego wykonuje zaangażowanie, jeżeli w terminie 30 dni nie zostanie zastąpiony innym stosunkiem obejmującym co najmniej minimalny zakres z §~\ref{sha:zaangazowanie} ust.~1;
    \item złożenia przez niego oświadczenia o zakończeniu zaangażowania;
    \item spełnienia przesłanki z §~\ref{sha:zaangazowanie} ust.~3;
    \item śmierci albo orzeczenia trwałej niezdolności do pracy.
  \end{enumerate}
  \item Ustanie zaangażowania stwierdza Rada Dyrektorów uchwałą podjętą bez udziału zainteresowanego, w terminie 30 dni od dnia, w którym dowiedziała się o zdarzeniu. Uchwała wskazuje Dzień Ustania Zaangażowania i liczbę Akcji Niezwolnionych obliczoną zgodnie z \US{9} ust.~2 i doręcza się ją Akcjonariuszowi Funkcjonalnemu pocztą elektroniczną.
  \item Odejście jest „Odejściem Uzasadnionym”, jeżeli zaangażowanie ustało wskutek śmierci, trwałej niezdolności do pracy, rozwiązania stosunku prawnego przez Spółkę bez istotnej przyczyny leżącej po stronie Akcjonariusza Funkcjonalnego albo odejścia uzgodnionego z Radą Dyrektorów. Pozostałe odejścia są „Odejściem Zwykłym”.
  \item Skutki Odejścia Uzasadnionego oraz śmierci i trwałej niezdolności do pracy określa \US{9} ust.~5: w razie śmierci albo trwałej niezdolności do pracy zwalniają się wszystkie Akcje Niezwolnione, a w pozostałych przypadkach Odejścia Uzasadnionego --- Akcje Niezwolnione, które zwolniłyby się w ciągu 12 miesięcy od Dnia Ustania Zaangażowania. Zwolnienie następuje z mocy Umowy Spółki; Założyciele zobowiązują się głosować za uchwałami potrzebnymi do jego ujawnienia.
  \item Akcjonariusz Funkcjonalny może zakwestionować uchwałę Rady Dyrektorów w terminie 14 dni od jej doręczenia, wnosząc o rozstrzygnięcie sporu w trybie §~\ref{sha:spory}. Do czasu rozstrzygnięcia ani Spółka, ani żaden dyrektor lub dalszy pełnomocnik nie wykonuje pełnomocnictwa, o którym mowa w §~\ref{sha:wykonanie} ust.~3, a wskazany nabywca nie przyjmuje oferty ani nie żąda zawarcia umowy zbycia co do akcji objętych sporem, chyba że do upływu terminu z \US{9} ust.~4 pozostało mniej niż 30 dni.
\end{enumerate}
\end{ContractClause}

\begin{ContractClause}{Wykonanie obowiązku zbycia akcji}
\label{sha:wykonanie}
\begin{enumerate}
  \item Akcjonariusz Funkcjonalny zobowiązuje się wobec Spółki do zbycia Akcji Niezwolnionych na rzecz nabywcy wskazanego zgodnie z \US{9} ust.~4, w terminach tam określonych (art.~393 Kodeksu cywilnego); nabywca może żądać zawarcia umowy zbycia bezpośrednio od Akcjonariusza Funkcjonalnego. W celu zabezpieczenia wykonania tego zobowiązania Akcjonariusz Funkcjonalny w dniu zawarcia niniejszej umowy, a w przypadku akcji serii~F --- w dniu zawarcia umowy objęcia akcji, składa ponadto nieodwołalną ofertę zbycia Akcji Objętych według wzoru z Załącznika A oraz udziela pełnomocnictwa, o którym mowa w ust.~3.
  \item Cena za jedną Akcję Niezwolnioną jest równa jej cenie emisyjnej. Strony potwierdzają, że cena ta wynika z Umowy Spółki i z mechanizmu stopniowego zwalniania akcji, przyjętych przy zawiązaniu Spółki jako warunek objęcia akcji przez wszystkich Założycieli, a nie z oceny wartości rynkowej akcji.
  \item Akcjonariusz Funkcjonalny udziela Spółce oraz każdemu dyrektorowi niewykonawczemu, z prawem do udzielania dalszych pełnomocnictw, pełnomocnictwa według wzoru z Załącznika B do zawarcia w jego imieniu umowy zbycia Akcji Niezwolnionych na rzecz nabywcy wskazanego zgodnie z \US{9} ust.~4 oraz do wydania dyspozycji wpisu nabywcy w rejestrze akcjonariuszy. Pełnomocnictwo jest nieodwołalne, ponieważ służy wykonaniu obowiązku wynikającego z Umowy Spółki i niniejszej umowy (art.~101 §~1 Kodeksu cywilnego), i nie wygasa ze śmiercią mocodawcy. Pełnomocnictwo może być wykonane wyłącznie po stwierdzeniu ustania zaangażowania zgodnie z §~\ref{sha:ustanie} i tylko co do Akcji Niezwolnionych.
  \item Nabywca zapłaci cenę w terminie 14 dni od zawarcia umowy zbycia na rachunek bankowy Akcjonariusza Funkcjonalnego wskazany w ofercie albo w późniejszym zawiadomieniu; dyspozycję wpisu nabywcy w rejestrze akcjonariuszy wydaje się dopiero po zapłacie. Złożenie ceny do depozytu sądowego jest dopuszczalne w przypadkach określonych w art.~467 Kodeksu cywilnego. Zbycie następuje w formie dokumentowej (art.~300$^{36}$ §~4 KSH), a nabycie --- z chwilą wpisu nabywcy do rejestru akcjonariuszy (art.~300$^{37}$ §~1 KSH).
  \item Spółka nie jest nabywcą Akcji Niezwolnionych ani nie finansuje ich nabycia. Jeżeli obowiązek zbycia wygaśnie zgodnie z \US{9} ust.~4, Spółka niezwłocznie zgłasza to do rejestru akcjonariuszy, a oferta i pełnomocnictwo w tym zakresie wygasają.
  \item Akcjonariusz Funkcjonalny nie obciąża Akcji Niezwolnionych i nie dokonuje czynności utrudniających wykonanie obowiązku zbycia.
\end{enumerate}
\end{ContractClause}

\begin{ContractClause}{Własność intelektualna}
\label{sha:ip}
\begin{enumerate}
  \item Każdy Założyciel, który wnosi do Spółki wkład niepieniężny w postaci praw własności intelektualnej, zawiera ze Spółką w organizacji, najpóźniej w dniu zawarcia Umowy Spółki, umowę przeniesienia tych praw według wzoru Umowy nr~1 z dokumentu „Umowy przeniesienia autorskich praw majątkowych --- Basilisk Systems P.S.A.” i przekazuje Spółce nośniki, repozytoria i dokumentację.
  \item Każdy Założyciel zawiera ze Spółką, najpóźniej w dniu rozpoczęcia zaangażowania, umowę przeniesienia praw do utworów tworzonych dla Spółki według wzoru Umowy nr~2 z tego dokumentu, chyba że prawa te przechodzą na Spółkę z mocy ustawy albo na podstawie umowy o pracę zawierającej postanowienia nie mniej korzystne dla Spółki.
  \item Założyciel oświadcza, że według jego najlepszej wiedzy: (a)~prawa wnoszone do Spółki nie przysługują jego obecnemu ani byłemu pracodawcy, uczelni ani zleceniodawcy; (b)~ujawnił Radzie Dyrektorów wszystkie znane mu prawa osób trzecich, w tym licencje oprogramowania otwartego wykorzystane w wnoszonych utworach; (c)~nie zobowiązał się wobec osób trzecich w sposób uniemożliwiający przeniesienie praw.
  \item Technologia Założyciela niewniesiona do Spółki podlega licencji i obowiązkom określonym w \US{14} ust.~2. Założyciel ujawnia ją Radzie Dyrektorów w formie pisemnej w terminie 30 dni od zawarcia niniejszej umowy.
  \item Założyciele współdziałają przy ustaleniu, które wkłady wpisuje się w Części~B1, a które w Części~B2 Załącznika nr~1 do Umowy Spółki, i każdy z nich ponosi sam skutki podatkowe swojego wkładu. Spółka udostępnia Założycielom dokumenty potrzebne do rozliczenia podatkowego i do złożenia wniosku o interpretację indywidualną.
\end{enumerate}
\end{ContractClause}

\begin{ContractClause}{Poufność}
\label{sha:poufnosc}
\begin{enumerate}
  \item Strony zachowują w poufności treść niniejszej umowy oraz informacje o Spółce, jej technologii, klientach, finansowaniu i akcjonariuszach, uzyskane w związku z niniejszą umową lub udziałem w Spółce.
  \item Obowiązek poufności nie dotyczy informacji: (a)~publicznie dostępnych bez naruszenia niniejszej umowy; (b)~ujawnianych doradcom związanym obowiązkiem poufności, potencjalnym inwestorom i nabywcom akcji, którzy zobowiązali się do poufności; (c)~których ujawnienia wymagają przepisy prawa albo organ władzy publicznej, po uprzednim powiadomieniu Spółki, jeżeli jest to dopuszczalne; (d)~zgłaszanych albo ujawnianych publicznie zgodnie z ustawą o ochronie sygnalistów, bez obowiązku uprzedniego powiadomienia Spółki; (e)~znanych Stronie zgodnie z prawem przed ich uzyskaniem w związku z udziałem w Spółce albo wytworzonych przez nią niezależnie.
  \item Obowiązek poufności trwa przez 5 lat od dnia, w którym Strona przestała być akcjonariuszem Spółki, a w odniesieniu do tajemnicy przedsiębiorstwa w rozumieniu ustawy o zwalczaniu nieuczciwej konkurencji --- dopóki informacja ma taki charakter.
\end{enumerate}
\end{ContractClause}

\begin{ContractClause}{Zakaz konkurencji i zakaz przejmowania}
\label{sha:konkurencja}
\begin{enumerate}
  \item Przez okres zaangażowania oraz --- w przypadku Odejścia Zwykłego --- przez 12 miesięcy od Dnia Ustania Zaangażowania Akcjonariusz Funkcjonalny nie prowadzi działalności konkurencyjnej wobec Spółki w zakresie przedmiotu działalności faktycznie wykonywanego przez Spółkę w Dniu Ustania Zaangażowania i nie uczestniczy w podmiocie prowadzącym taką działalność, z wyjątkiem inwestycji portfelowych nieprzekraczających 5\% udziału.
  \item Przez okres zaangażowania oraz przez 12 miesięcy od Dnia Ustania Zaangażowania, niezależnie od rodzaju odejścia, Akcjonariusz Funkcjonalny nie nakłania aktywnie pracowników ani kluczowych współpracowników Spółki do zakończenia współpracy ze Spółką ani klientów i partnerów Spółki do jej ograniczenia. Zatrudnienie osoby, która sama odpowiedziała na publiczne ogłoszenie, nie narusza tego zakazu.
  \item Rada Dyrektorów może, uchwałą podjętą bez udziału zainteresowanego, zwolnić Akcjonariusza Funkcjonalnego z zakazu w całości albo w części.
  \item W zakresie, w jakim Akcjonariusz Funkcjonalny jest pracownikiem Spółki, zakaz konkurencji po ustaniu stosunku pracy wymaga odrębnej umowy z odszkodowaniem zgodnie z art.~101$^{2}$ Kodeksu pracy; niniejszy paragraf obowiązuje wówczas w zakresie, w jakim nie jest z nią sprzeczny.
\end{enumerate}
\end{ContractClause}

\begin{ContractClause}{Kary umowne}
\label{sha:kary}
\begin{enumerate}
  \item Strona zapłaci Spółce karę umowną: (a)~25\,000 zł za każde naruszenie §~\ref{sha:poufnosc}; (b)~50\,000 zł za każde naruszenie §~\ref{sha:konkurencja}; (c)~25\,000 zł za każde naruszenie §~\ref{sha:wykonanie} ust.~6; w razie naruszenia ciągłego --- dodatkowo 5\,000 zł za każdy pełny miesiąc jego trwania. Łączna wysokość kar należnych od jednej Strony nie przekracza 150\,000 zł.
  \item Kwoty określone w ust.~1 podlegają waloryzacji z dniem 1 marca każdego roku, począwszy od roku następującego po roku zawarcia umowy, o średnioroczny wskaźnik cen towarów i usług konsumpcyjnych ogółem za poprzedni rok kalendarzowy, ogłaszany przez Prezesa Głównego Urzędu Statystycznego w Dzienniku Urzędowym Rzeczypospolitej Polskiej „Monitor Polski”, jeżeli wskaźnik ten przekracza 100. Zwaloryzowane kwoty zaokrągla się w górę do pełnych 100 zł. Wysokość kary ustala się według kwot obowiązujących w dniu naruszenia, a w razie naruszenia ciągłego --- w pierwszym dniu każdego miesiąca jego trwania. Waloryzacja nie wymaga zmiany umowy.
  \item Zapłata kary nie wyłącza dochodzenia odszkodowania przenoszącego jej wysokość ani roszczeń o zaniechanie naruszeń.
\end{enumerate}
\end{ContractClause}

\begin{ContractClause}{Finansowanie i kolejne rundy}
\label{sha:finansowanie}
\begin{enumerate}
  \item Założyciele dążą do finansowania Spółki w pierwszej kolejności z grantów, przychodów i emisji w granicach limitów z \US{5}. Żaden Założyciel nie jest zobowiązany do dopłat ani do udzielania Spółce pożyczek.
  \item Pożyczka udzielona Spółce przez Założyciela wymaga uchwały Rady Dyrektorów i jest udzielana na warunkach rynkowych w formie pisemnej. Jeżeli pożyczkodawca jest dyrektorem, umowa wymaga ponadto zgody Walnego Zgromadzenia (\US{13} ust.~1 lit.~j), a po stronie Spółki zawiera ją pełnomocnik powołany uchwałą akcjonariuszy (art.~300$^{79}$ §~1 KSH, \US{11} ust.~3).
  \item Przed wprowadzeniem inwestora Założyciele uzgadniają projekt zmiany Umowy Spółki ustanawiającej serię inwestorską (\US{5} ust.~4) oraz zmiany niniejszej umowy wymagane przez inwestora. Żaden Założyciel nie jest zobowiązany do zgody na postanowienia pogarszające jego sytuację w stosunku do innych Założycieli tej samej serii.
\end{enumerate}
\end{ContractClause}

\begin{ContractClause}{Spory i impas}
\label{sha:spory}
\begin{enumerate}
  \item Strony rozwiązują spory w dobrej wierze w drodze negocjacji. Jeżeli spór nie zostanie rozwiązany w terminie 30 dni od wezwania, Strony poddają go mediacji prowadzonej przez \field{ośrodek mediacyjny}.
  \item Jeżeli Walne Zgromadzenie dwukrotnie, w odstępie co najmniej 30 dni, nie podejmie uchwały w sprawie zastrzeżonej z \US{13} ust.~1 lit.~b albo~h, a sprawa jest niezbędna dla dalszego prowadzenia działalności Spółki, każda ze Stron może wszcząć mediację z ust.~1; Strony uczestniczą w niej w dobrej wierze.
  \item W razie bezskuteczności mediacji spór rozstrzyga sąd powszechny właściwy dla siedziby Spółki. Postanowienie to nie ogranicza prawa do żądania zabezpieczenia.
\end{enumerate}
\end{ContractClause}

\begin{ContractClause}{Czas trwania}
\label{sha:czas}
\begin{enumerate}
  \item Umowa zostaje zawarta na czas oznaczony 10 lat od dnia jej zawarcia i przedłuża się na kolejne okresy 5-letnie, chyba że Strona najpóźniej na 12 miesięcy przed upływem okresu oświadczy pozostałym Stronom, że nie przedłuża umowy w stosunku do siebie. Umowa wiąże Stronę, dopóki jest ona akcjonariuszem Spółki; §~\ref{sha:poufnosc}--\ref{sha:kary} i \ref{sha:spory} wiążą ją także po tej chwili w terminach w nich określonych.
  \item Umowa wygasa w całości z chwilą zbycia wszystkich akcji Spółki na rzecz jednego nabywcy albo grupy nabywców działających w porozumieniu, połączenia Spółki z inną spółką, w której Spółka nie jest przejmującą, albo zakończenia likwidacji Spółki.
  \item Przed upływem okresu, o którym mowa w ust.~1, Strona może wypowiedzieć umowę w stosunku do siebie wyłącznie z ważnych powodów, ze skutkiem na koniec miesiąca następującego po doręczeniu wypowiedzenia. Wypowiedzenie albo nieprzedłużenie umowy nie zwalnia Strony z obowiązków związanych z akcjami wynikających z Umowy Spółki ani z obowiązków z §~\ref{sha:poufnosc}--\ref{sha:kary}.
  \item W razie wygaśnięcia, wypowiedzenia albo nieprzedłużenia umowy w stosunku do Strony zachowują moc: rozliczenia między Stronami, obowiązki dotyczące przeniesienia praw do utworów powstałych wcześniej (§~\ref{sha:ip}), obowiązki związane z ustaniem zaangażowania, które nastąpiło wcześniej (§~\ref{sha:ustanie} i~\ref{sha:wykonanie}), wraz z udzielonymi w tym celu ofertą i pełnomocnictwem, oraz powstałe wcześniej roszczenia.
\end{enumerate}
\end{ContractClause}

\begin{ContractClause}{Postanowienia końcowe}
\label{sha:koncowe}
\begin{enumerate}
  \item Zmiana niniejszej umowy wymaga formy pisemnej pod rygorem nieważności i zgody wszystkich Stron. Zmiana załączników oraz dostosowanie odesłań do zmiany Umowy Spółki wymaga zgody Stron reprezentujących co najmniej 75\% wszystkich głosów w Spółce, jeżeli nie zwiększa obowiązków ani nie ogranicza praw żadnej ze Stron.
  \item Zawiadomienia dokonuje się pocztą elektroniczną na adresy wskazane w komparycji albo w oświadczeniu o przystąpieniu, a w braku adresu --- przesyłką poleconą. Zmiana adresu nie stanowi zmiany umowy i jest skuteczna od chwili zawiadomienia pozostałych Stron.
  \item Prawem właściwym jest prawo polskie. Jeżeli postanowienie umowy jest nieważne albo niewykonalne, pozostałe zachowują moc, a Strony zastąpią je postanowieniem najbliższym jego celowi gospodarczemu.
  \item Umowę sporządzono w \field{liczba} jednobrzmiących egzemplarzach, po jednym dla każdej Strony; dopuszcza się zawarcie umowy przy użyciu kwalifikowanych podpisów elektronicznych. Załączniki A--D stanowią jej integralną część.
  \item Spółka jest administratorem danych osobowych Stron i przetwarza je w celu wykonania niniejszej umowy oraz prowadzenia dokumentacji korporacyjnej (art.~6 ust.~1 lit.~b i~c RODO); pozostałe informacje wymagane przez art.~13 RODO Spółka przekazuje Stronom przy zawarciu umowy.
  \item Umowę sporządzono w języku polskim i angielskim; w razie rozbieżności rozstrzyga wersja polska.
  \item Założyciel pozostający we wspólności majątkowej małżeńskiej składa przy zawarciu niniejszej umowy oświadczenie o przynależności swoich akcji według Części~1 Załącznika D. Akcje objęte w zamian za prawa twórcy należą co do zasady do jego majątku osobistego (art.~33 pkt~9--10 Kodeksu rodzinnego i opiekuńczego). Jeżeli część akcji należy do majątku wspólnego, Założyciel przedstawia ponadto zgodę małżonka według Części~2 Załącznika D; to samo dotyczy osoby przystępującej do umowy.
\end{enumerate}
\end{ContractClause}

\vspace{1cm}
\begin{center}
\begin{tabularx}{\textwidth}{XX}
\field{podpis Założyciela A} & \field{podpis Założyciela B} \\[1.2cm]
\field{podpis Założyciela C} & \field{podpis Założyciela D} \\[1.2cm]
\field{podpis Założyciela E} & \field{za Spółkę: pełnomocnik} \\
\end{tabularx}
\end{center}

\newpage
\section*{Załącznik A --- Wzór nieodwołalnej oferty zbycia akcji}
\addcontentsline{toc}{section}{Załącznik A --- Wzór nieodwołalnej oferty zbycia akcji}

\field{miejscowość, data}

Ja, \field{imię, nazwisko, PESEL, adres}, akcjonariusz \CompanyFull{} („Spółka”), działając na podstawie \US{9} oraz §~\ref{sha:wykonanie} umowy akcjonariuszy z dnia \field{data}, niniejszym składam nieodwołalną ofertę zbycia \field{liczba} akcji Spółki serii~\field{seria} o numerach \field{od--do} („Akcje Objęte”) na następujących warunkach:

\begin{enumerate}
  \item Oferta jest skierowana do nabywcy wskazanego przez Radę Dyrektorów Spółki, a w braku wskazania --- do Założycieli Pierwotnych, zgodnie z \US{9} ust.~4.
  \item Oferta może zostać przyjęta wyłącznie co do Akcji Niezwolnionych w rozumieniu umowy akcjonariuszy, po stwierdzeniu ustania mojego zaangażowania, w terminach z \US{9} ust.~4.
  \item Cena za jedną akcję jest równa jej cenie emisyjnej i wynosi \field{kwota} zł; cenę należy zapłacić w terminie 14 dni od przyjęcia oferty na rachunek nr \field{numer rachunku}.
  \item Oferta wiąże mnie i moich następców prawnych do dnia, w którym wszystkie Akcje Objęte zostaną zwolnione z obowiązku zbycia albo obowiązek zbycia wygaśnie.
\end{enumerate}

\field{podpis}

\newpage
\section*{Załącznik B --- Wzór pełnomocnictwa}
\addcontentsline{toc}{section}{Załącznik B --- Wzór pełnomocnictwa}

\field{miejscowość, data}

Ja, \field{imię, nazwisko, PESEL, adres}, udzielam \CompanyFull{} oraz każdemu z jej dyrektorów niewykonawczych, z prawem udzielania dalszych pełnomocnictw, pełnomocnictwa do:

\begin{enumerate}
  \item zawarcia w moim imieniu umowy zbycia Akcji Niezwolnionych, w liczbie stwierdzonej uchwałą Rady Dyrektorów podjętą zgodnie z §~\ref{sha:ustanie} umowy akcjonariuszy z dnia \field{data}, na rzecz nabywcy wskazanego zgodnie z \US{9} ust.~4, po cenie równej cenie emisyjnej;
  \item wydania dyspozycji wpisu nabywcy w rejestrze akcjonariuszy oraz dokonania innych czynności niezbędnych do przeniesienia tych akcji.
\end{enumerate}

Pełnomocnictwo jest nieodwołalne, ponieważ służy wykonaniu obowiązku wynikającego z umowy spółki i umowy akcjonariuszy (art.~101 §~1 Kodeksu cywilnego), nie wygasa z moją śmiercią i może być wykonane wyłącznie w przypadku i w trybie określonym w §~\ref{sha:wykonanie} umowy akcjonariuszy. Pełnomocnik może być drugą stroną umowy zbycia albo reprezentować obie strony.

\field{podpis}

\newpage
\section*{Załącznik C --- Wzór oświadczenia o przystąpieniu}
\addcontentsline{toc}{section}{Załącznik C --- Wzór oświadczenia o przystąpieniu}

\field{miejscowość, data}

Ja, \field{imię, nazwisko albo firma, PESEL albo KRS, adres, adres e-mail}, w związku z \field{nabyciem / objęciem} \field{liczba} akcji \CompanyFull{} serii~\field{seria} o numerach \field{od--do}, oświadczam, że przystępuję do umowy akcjonariuszy z dnia \field{data}, której treść otrzymałem, i przyjmuję wynikające z niej prawa i obowiązki w całości.

Jeżeli nabywane akcje są Akcjami Objętymi, składam jednocześnie ofertę i udzielam pełnomocnictwa według wzorów z Załączników A i~B.

\field{podpis}

Przyjmuję przystąpienie w imieniu Stron: \field{za Spółkę}


\newpage
\section*{Załącznik D --- Wzór oświadczenia i zgody małżonka}
\addcontentsline{toc}{section}{Załącznik D --- Wzór oświadczenia i zgody małżonka}

\field{miejscowość, data}

\textbf{Część 1. Oświadczenie Założyciela o przynależności akcji} (składa każdy Założyciel pozostający we wspólności majątkowej małżeńskiej):

Ja, \field{imię, nazwisko, PESEL}, oświadczam, że akcje \CompanyFull{} o numerach \field{od--do}, objęte w zamian za \field{prawa twórcy (aport IP) / środki z majątku osobistego}, należą do mojego majątku osobistego (art.~33 pkt~9--10 Kodeksu rodzinnego i opiekuńczego), a akcje o numerach \field{od--do}, objęte \field{za środki z majątku wspólnego}, należą do majątku wspólnego.

\field{podpis Założyciela}

\textbf{Część 2. Zgoda małżonka} (wyłącznie gdy część akcji należy do majątku wspólnego):

Ja, \field{imię, nazwisko, PESEL, adres}, małżonek \field{imię, nazwisko Założyciela}, w odniesieniu do akcji należących do majątku wspólnego, o numerach \field{od--do}, oświadczam, że:

\begin{enumerate}
  \item znam treść Umowy Spółki \CompanyFull{} oraz umowy akcjonariuszy z dnia \field{data};
  \item wyrażam zgodę na zawarcie umowy akcjonariuszy przez mojego małżonka w zakresie tych akcji, w tym na złożenie oferty zbycia i udzielenie pełnomocnictwa według wzorów z Załączników A i~B oraz na zbycie Akcji Niezwolnionych po cenie emisyjnej na warunkach tej umowy i Umowy Spółki;
  \item zgoda nie oznacza mojego przystąpienia do umowy akcjonariuszy ani przyjęcia przeze mnie zobowiązań pieniężnych z niej wynikających; jeżeli w wyniku podziału majątku wspólnego nabędę akcje, przystąpienie do umowy wymaga odrębnego oświadczenia według wzoru z Załącznika~C.
\end{enumerate}

\field{podpis małżonka}

\end{document}
